% This file is part of the modular source for Stratified Relational Workspace Theory.
% !TeX root = ../Stratified_Relational_Workspace_Theory.tex
\theorypart{X}{Persistence}

The preceding parts identify exact and approximate workspace subspaces at a
single time. Persistence adds three further requirements. The structural
subspace must remain trackable while the realization changes, actual states
must remain near that moving subspace, and the transported content must not
collapse to a trivial zero state. These are distinct conditions.

\topic{32}{Coherence-Error Dynamics}

\subtopic{32.1 Time-indexed transfer systems}

At time \(t\), let
\[
\mathcal C_V(t)=\bigoplus_{i\in V}C_i(t),
\qquad
\mathcal C_E(t)=\bigoplus_{e\in E}C_{t(e)}(t),
\]
and let
\[
D_t:\mathcal C_V(t)\longrightarrow\mathcal C_E(t),
\qquad
L_t=D_t^*D_t
\]
be the transfer discrepancy and transfer Laplacian. At time \(t\), the
combinatorial graph is held fixed while the local representations, transfers,
and metrics may vary. Changes of graph are treated later by a graph
correspondence.

To compare successive times, choose unitary identification maps
\[
U_t^V:\mathcal C_V(t)\longrightarrow\mathcal C_V(t+1),
\qquad
U_t^E:\mathcal C_E(t)\longrightarrow\mathcal C_E(t+1).
\]
They are comparison maps, not assumptions that the physical carriers are
identical. Define the one-step diagram defect
\[
\boxed{
\Delta_t^D
=
D_{t+1}U_t^V-U_t^ED_t,
\qquad
\delta_t^D=\|\Delta_t^D\|.
}
\]
Exact temporal compatibility is the condition \(\Delta_t^D=0\).

\begin{proposition}[Temporal Laplacian drift]\label{prop:temporal-laplacian-drift}
The carried Laplacian satisfies
\[
\boxed{
\left\|
(U_t^V)^*L_{t+1}U_t^V-L_t
\right\|
\leq
2\|D_t\|\delta_t^D+(\delta_t^D)^2.
}
\]
\end{proposition}

\begin{proof}
Set
\[
B_t=(U_t^E)^*\Delta_t^D.
\]
Because the comparison maps are unitary,
\[
(U_t^E)^*D_{t+1}U_t^V=D_t+B_t,
\qquad
\|B_t\|=\delta_t^D.
\]
Hence
\[
(U_t^V)^*L_{t+1}U_t^V
=(D_t+B_t)^*(D_t+B_t).
\]
Subtract \(D_t^*D_t\) and apply the triangle inequality.
\end{proof}

Let \(P_t\) project onto a selected rank-\(d\) low-energy workspace cluster
of \(L_t\), and set
\[
\mathcal W_t=\ran P_t.
\]
Write
\[
g_t=\lambda_{d+1}(L_t)-\lambda_d(L_t)>0
\]
for the separating spectral gap and
\[
\eta_t^L=2\|D_t\|\delta_t^D+(\delta_t^D)^2.
\]

\begin{corollary}[Temporal workspace tracking]\label{cor:temporal-workspace-tracking}
If
\[
\eta_t^L<\frac{g_t}{2},
\]
then the transported and newly extracted workspace projectors satisfy
\[
\boxed{
q_t^{\mathrm{str}}
:=
\left\|
P_{t+1}-U_t^VP_t(U_t^V)^*
\right\|
\leq
\frac{\eta_t^L}{g_t-\eta_t^L}
\leq
\frac{2\eta_t^L}{g_t}.
}
\]
\end{corollary}

\begin{proof}
Carry \(L_{t+1}\) back to time \(t\) using \(U_t^V\), apply
Theorem~\ref{thm:spectral-stability}, and conjugate the resulting projector
bound forward.
\end{proof}

The quantity \(q_t^{\mathrm{str}}\) measures structural motion after the
declared coordinate change has been removed. A workspace may therefore move
substantially in raw coordinates while remaining structurally unchanged.

\subtopic{32.2 State tracking inside a moving workspace}

Let the realized vertex state satisfy
\[
x_{t+1}=F_tx_t+\xi_t,
\]
where \(F_t:\mathcal C_V(t)\to\mathcal C_V(t+1)\) is the actual update and
\(\xi_t\) is an unmodeled disturbance. Define the off-workspace error
\[
\boxed{
e_t=\|(I-P_t)x_t\|.
}
\]
Introduce
\[
\alpha_t
=
\left\|
(I-P_{t+1})F_t(I-P_t)
\right\|,
\]
which measures persistence of prior off-workspace error, and
\[
\ell_t
=
\left\|
(I-P_{t+1})F_tP_t
\right\|,
\]
which measures leakage of a valid workspace state under the update.

\begin{proposition}[Grounded coherence recurrence]\label{prop:grounded-coherence-recurrence}
Every trajectory satisfies
\[
\boxed{
e_{t+1}\leq\alpha_te_t+\sigma_t,}
\]
where
\[
\boxed{
\sigma_t
=
\ell_t\|P_tx_t\|
+
\|(I-P_{t+1})\xi_t\|.
}
\]
\end{proposition}

\begin{proof}
Insert \(I=P_t+(I-P_t)\) before \(x_t\), apply \(I-P_{t+1}\), and use
submultiplicativity:
\[
\begin{aligned}
e_{t+1}
&\leq
\|(I-P_{t+1})F_t(I-P_t)x_t\|\\
&\quad+
\|(I-P_{t+1})F_tP_tx_t\|
+
\|(I-P_{t+1})\xi_t\|.
\end{aligned}
\]
The displayed recurrence follows from the definitions of \(\alpha_t\) and
\(\ell_t\).
\end{proof}

This derivation gives the contraction and disturbance terms an operational
meaning. The scalar recurrence is not an independent postulate. It follows
from an update rule, a moving workspace projector, and a declared disturbance
model.

\subtopic{32.3 Persistence theorem \(T_6\)}

\begin{theorem}[Nonautonomous persistence theorem]\label{thm:persistence}
Let \(e_t\geq0\) satisfy
\[
e_{t+1}\leq\alpha_te_t+\sigma_t,
\qquad
\alpha_t,\sigma_t\geq0.
\]
Then, for every \(t\geq1\),
\[
\boxed{
\begin{aligned}
e_t
\leq
&\left(\prod_{r=0}^{t-1}\alpha_r\right)e_0\\
&+
\sum_{j=0}^{t-1}
\left(\prod_{r=j+1}^{t-1}\alpha_r\right)\sigma_j,
\end{aligned}
}
\]
where an empty product equals one. If
\[
\alpha_t\leq\bar\alpha<1,
\qquad
\sigma_t\leq\bar\sigma,
\]
then
\[
\boxed{
e_t
\leq
\bar\alpha^te_0
+
\frac{\bar\sigma(1-\bar\alpha^t)}{1-\bar\alpha},
}
\]
and consequently
\[
\boxed{
\limsup_{t\to\infty}e_t
\leq
\frac{\bar\sigma}{1-\bar\alpha}.
}
\]
\end{theorem}

\begin{proof}
Repeated substitution yields
\[
\begin{aligned}
e_t
&\leq
\alpha_{t-1}\cdots\alpha_0e_0
+
\alpha_{t-1}\cdots\alpha_1\sigma_0
+
\cdots
+
\sigma_{t-1},
\end{aligned}
\]
which is the first formula. Under the uniform bounds, dominate each product by
the corresponding power of \(\bar\alpha\) and sum the geometric series.
\end{proof}

\begin{corollary}[Vanishing-disturbance convergence]\label{cor:vanishing-disturbance}
If \(\alpha_t\leq\bar\alpha<1\) and \(\sigma_t\to0\), then
\[
e_t\to0.
\]
More generally, if \(\sum_{t\geq0}\sigma_t<\infty\), then the disturbance
contribution in Theorem~\ref{thm:persistence} vanishes asymptotically.
\end{corollary}

\begin{proof}
Split the convolution sum into a finite early part, whose coefficients decay
to zero, and a late tail on which \(\sigma_t\) is uniformly small. The
summable case follows by the same argument.
\end{proof}

\subtopic{32.4 Mean-square persistence}

For stochastic systems, pathwise bounded disturbances may be too strong. Let
\(\{\mathcal F_t\}\) be a filtration containing the history through time
\(t\).

\begin{proposition}[Mean-square persistence]\label{prop:mean-square-persistence}
Suppose
\[
\mathbb E[e_{t+1}^2\mid\mathcal F_t]
\leq
q e_t^2+v_t,
\qquad
0\leq q<1,
\]
and \(\mathbb E[v_t]\leq\bar v\). Then
\[
\boxed{
\mathbb E[e_t^2]
\leq
q^t\mathbb E[e_0^2]
+
\frac{\bar v(1-q^t)}{1-q},
}
\]
so
\[
\limsup_{t\to\infty}\mathbb E[e_t^2]
\leq
\frac{\bar v}{1-q}.
\]
\end{proposition}

\begin{proof}
Take expectations, iterate the resulting deterministic recurrence, and sum a
geometric series.
\end{proof}

\subtopic{32.5 The non-collapse condition}

Low off-workspace error is not sufficient for meaningful persistence. The
zero update \(F_t=0\) drives every state exactly into every linear workspace
while destroying all content. Define the internal transport
\[
B_t=P_{t+1}F_tP_t:\mathcal W_t\longrightarrow\mathcal W_{t+1}
\]
and its non-collapse margin
\[
\boxed{
\mu_t^{\mathrm{keep}}
=
\sigma_{\min}(B_t|_{\mathcal W_t}).
}
\]
A content-preserving persistence claim requires
\[
\mu_t^{\mathrm{keep}}>0
\]
for the relevant interval, or an information-theoretic replacement when the
workspace state is stochastic. Structural tracking, state tracking, and
content retention are therefore separate gates.

A finite-horizon persistence certificate may be recorded as
\[
\boxed{
\mathcal P_{0:T}
=
\left(
\max_{t<T}q_t^{\mathrm{str}},
\max_{t\leq T}e_t,
\min_{t<T}\mu_t^{\mathrm{keep}},
\sum_{t<T}\sigma_t
\right).
}
\]
No single coordinate of this vector subsumes the others.

\topic{33}{Interpretation of the Contraction Parameter}

The inequality
\[
\alpha_t<1
\]
means that deviations already outside the current workspace are damped by the
next update. Mechanisms that can reduce \(\alpha_t\) or \(\sigma_t\) include
\[
\begin{aligned}
&\text{verification and error feedback},\\
&\text{shared notation and stable interfaces},\\
&\text{artifact persistence and version control},\\
&\text{proof checking and experimental replication},\\
&\text{handoff reconstruction and redundancy}.
\end{aligned}
\]
Sources of disturbance include
\[
\begin{aligned}
&\text{context loss and representational drift},\\
&\text{miscommunication and transfer-map error},\\
&\text{participant or model replacement},\\
&\text{artifact corruption},\\
&\text{unsupported inference and objective drift}.
\end{aligned}
\]

\subtopic{33.1 Four persistence levels}

The theory distinguishes:
\[
\begin{array}{ll}
\text{structural persistence}
& q_t^{\mathrm{str}}\text{ remains small},\\
\text{state persistence}
& e_t\text{ remains small},\\
\text{content persistence}
& \mu_t^{\mathrm{keep}}\text{ remains bounded away from zero},\\
\text{operational persistence}
& \text{predictive and causal capacities remain matched}.
\end{array}
\]
The first three are developed in this part. Operational matching is incorporated
into the identity theory of Part~XII.

\subtopic{33.2 Moving-line example}

Let \(\mathcal C_V(t)=\mathbb R^2\) and let \(P_t\) project onto
\[
u_t=(\cos\theta_t,\sin\theta_t).
\]
If \(U_t^V\) is rotation through
\(\theta_{t+1}-\theta_t\) and \(F_t=U_t^V\), then a state
\(x_t=u_t\) obeys
\[
e_t=e_{t+1}=0,
\qquad
\mu_t^{\mathrm{keep}}=1.
\]
The workspace changes raw coordinates but not abstract identity. If instead
\(F_t=I\), then
\[
e_{t+1}
=
|\sin(\theta_{t+1}-\theta_t)|.
\]
The same structural motion now produces a state-tracking defect because the
actual update fails to transport the content with the workspace.
